\head{EGA V 预备笔记：超平面截面与锥投影}
\source{\texttt{EGAVaiello.pdf} 第1--12页；本篇与 \texttt{EGAV5.ps} 的开头部分对应}
\noindent\textbf{底本标题署名：}Alexandre Grothendieck 的 EGA V；其“预备笔记”的翻译与编辑：Piotr Blass、Joseph Blass。

\egasection{编写摘要}
这一稿把应收入最终稿的全部结果杂陈在一起，作了详细概述。要达到最终稿的形式，需要对目前这个“第零阶段”彻底重新组织。第一步大概应当是重订计划（毫无疑问，现在的第11、12、14、15节将大幅提前）。我甚至还没有写第16节；原则上，它应当没有困难，也不以任何方式影响前面的各节，因为涉及的只是转换表述。

你会注意到，这里有一条命题10.3，它应当出现在前面的某个段落。

就此，我还想告诉你，我另有若干彼此很不相同、但都涉及双有理映射的结果，很愿意把它们放到某处。我觉得这些还不足以独立构成一节。你对放在哪里有什么建议吗？我打算不久把它们连同这份笔记的第16节寄给你。\ednote{询问 AG：第16节究竟是否写过。}

此外，目前的第20段大概会扩展成两段，其中一段由 Grassmann 概形上属于“初等几何”类型的结果组成。必要时，是否也可以把我跟你提过的关于双有理变换的补充材料放在这里（如果找不到更合适的地方）？

\egasection{拟定目录}
\begin{enumerate}[label=\arabic*)]
\item 记号。
\item 泛超平面截面：局部性质。
\item 泛超平面截面：几何不可约性与连通性。
\item 变动超平面截面：“充分一般”的截面。
\item Seidenberg 类型定理。
\item 任意超平面截面的连通性。
\item 特殊超平面截面与多截面构造的应用。
\item 例外超平面集合的维数。
\item 射影浸入的改变。
\item 超平面截面的铅笔与爆破簇的纤维化。
\item Grassmann 概形。
\item 前述结果向线性截面的推广。
\item 初等态射与 M. Artin 的一个定理。
\item 锥投影。
\item 部分前述结果的公理化。
\item 转换为线性系的语言。
\end{enumerate}
以上是底本列出的拟定目录，不表示全部节均已经写成，也不表示本批已翻译目录中全部内容。

\egasection{§1\quad 预备知识与记号}
设 $S$ 为预概形，$E$ 为 $S$ 上有限型局部自由模，$E^\vee$ 为其对偶。记 $P=\mathbf P(E)$ 为 $E$ 定义的射影丛，$P^\vee$ 为 $E^\vee$ 定义的射影丛。我们称 $P^\vee$ 为 $P$ 的\textbf{超平面概形}。这个术语可作如下解释。

设 $\xi$ 是 $P^\vee$ 在 $S$ 上的截面，因而由 $E^\vee$ 的一个可逆商模 $L$ 确定。由它得到 $(E^\vee)_P=(E_P)^\vee$ 的可逆商模 $L_P$；另一方面，我们有 $E_P$ 的可逆商模 $\OO_P(1)$。取对偶后，可以把 $L_P^{-1}$、$\OO_P(-1)$ 分别看成 $E_P$、$(E_P)^\vee$ 的可逆子模，且局部为直和因子。配对 $E_P\otimes(E_P)^\vee\to\OO_P$ 因而定义自然配对
\begin{equation}
 \OO_P(-1)\otimes L_P^{-1}\longrightarrow\OO_P,\tag{*}
\end{equation}
或者其转置同态
\begin{equation}
 \OO_P\longrightarrow\OO_P(1)\otimes L_P=L_P(1).\tag{**}
\end{equation}
后者就是一个由 $\xi$ 典范确定的 $L_P(1)$ 的截面。此截面的“除子”，即由（*）的像理想定义的 $P$ 的闭子概形 $H_\xi$，称为点 $\xi\in P^\vee(S)$ 定义的\textbf{超平面}。

还可以这样描述它：在 $S$ 上局部地，$\xi$ 由 $E$ 的一个截面 $\phi$ 给出，且对每个 $s$ 都有 $\phi(s)\ne0$；$\xi$ 确定 $\phi$，相差一个 $\OO_S$ 可逆截面的倍数。由于 $E=p_*\OO_P(1)$（$p:P\to S$ 为投影），$\phi$ 也可视为 $\OO_P(1)$ 的截面，而它的除子正是 $H_\xi$。

当然，若把 $L^{-1}$ 看成 $E$ 的可逆子模，且局部为直和因子，则 $\xi$（即 $L$ 或 $L^{-1}\subset E$）与 $\phi$ 的对应，由取 $L^{-1}$ 的一个处处不为零的截面得到，也就是选择 $L^{-1}$ 的一个平凡化；这种平凡化无论如何总在局部存在。注意，$H_\xi$ 也就是 $\mathbf P(E/L^{-1})$（典范同构）。这是第三种描述方式。特别地，$\mathbf P(E/L^{-1})$ 典范嵌入 $P=\mathbf P(E)$，这一描述还证明了 $H_\xi$ 是 $S$ 上的射影丛，因而在 $S$ 上光滑。（还应当在 EGA IV 第17段中说一下，射影丛是光滑的。）无疑，从这一描述开始会更好。

\noindent\textbf{备注。}
$H_\xi$ 的构造与基变换相容；换言之，得到函子范畴 $(\mathrm{Sch}/S)^{\mathrm{op}}\to\mathrm{Ens}$ 中的自然变换
\[
 P^\vee\longrightarrow\operatorname{Div}(P/S),
\]
其中右端表示 $P/S$ 的“相对除子”函子：它在任意 $S$-预概形 $S'$ 上的值，是 $P_{S'}$ 中相对 $S'$ 为横截完全交且余维为1的闭子概形集合（比较第19段）。\ednote{这里使用 EGA I 新版的记号；第19段指 EGA IV。}
容易证明这是函子的单态射，也就是说，$\xi$ 由 $H_\xi$ 确定。这说明“超平面概形”的术语是恰当的。我们将看到，$\operatorname{Div}(P/S)$ 可由各 $\mathbf P(\operatorname{Sym}^m_k(E^\vee))$ 的不交并预概形表示，从而 $P^\vee$ 可视为其一个既开又闭的子概形。\ednote{比较 Mumford 的 \textit{Lectures on Curves on an Algebraic Surface}。}\cmt{P03}{底本此处在一般底预概形 $S$ 上使用了未定义的下标 $k$，对称幂的对偶位置和秩的退化情形也需核对。这段是后文拟证明的宣告，不能仅据此视为已完成证明。本文按所上传英文稿保留表达。}
（严格说，对 $P/S$ 相对除子的确定，用目前已有的工具和第II章的结果就能进行，也可以作为例子补到 EGA IV 第19段中。）

现在作基变换 $S'=P^\vee\to S$，并取 $P^\vee_{S'}=\mathbf P(E^\vee_{S'})$ 在 $S'$ 上的对角截面，或称“泛截面”。得到 $P_{S'}=P\times_SP^\vee$ 的闭子概形 $H$，有时称为 $P$ 与 $P^\vee$ 的\textbf{关联概形}。它由典范同态
\[
 \OO_P(-1)\otimes_S\OO_{P^\vee}(-1)
 \longrightarrow\OO_{P\times_SP^\vee}
\]
的像理想定义。因此它是 $P^\vee$ 上的射影丛，由对称性，也是 $P$ 上的射影丛。

注意，给定 $P^\vee$ 在 $S$ 上的截面 $\xi$ 后，可以从这个“普遍超平面”$H$ 出发，沿基变换 $\xi:S\to P^\vee$ 取逆像，恢复“特殊”超平面 $H_\xi$。对 $P^\vee$ 的每个取值于任意 $S$-预概形 $S'$ 的点也一样：把它看成 $P^\vee_{S'}$ 在 $S'$ 上的截面，便定义 $H_\xi\subset P_{S'}$；它正是沿 $\xi:S'\to P^\vee$ 对 $H$ 作基变换的逆像。

以下假设有一个在 $P$ 上有限型的预概形 $X$，以及一个 $S$-态射 $f:X\to P$。\ednote{或者应当是“在 $S$ 上有限型”，我不能确定。}
这一段的主要目标之一，是对每个超平面 $H_\xi$ 研究其逆像
\[
 Y_\xi=f^{-1}(H_\xi)=X\times_PH_\xi,
\]
特别是联系 $X$ 与 $Y_\xi$ 的性质。像通常一样，也须考虑取值于任意 $S$-概形 $S'$ 的点；此时 $H_\xi$ 为 $P_{S'}$ 中的超平面，并置
\[
 Y_\xi=f_{S'}^{-1}(H_\xi)=X_{S'}\times_{P_{S'}}H_\xi
       =X\times_PH_\xi.
\]
下标 $S'$ 表示基变换 $S'\to S$ 的效果；在最后一式中，通过复合 $H_\xi\to P_{S'}\to P$ 把 $H_\xi$ 视为 $P$-概形。因而仍宜先考虑 $\xi$“普遍”的情况：$S'=P^\vee$，$\xi$ 为对角截面，于是 $H_\xi=H$；此时记 $Y=Y_\xi$（等待 Dieudonné 提出更合适的记号）。对于一般的 $\xi:S'\to P^\vee$，也有 $Y_\xi=Y\times_{P^\vee}S'$。

最后，若 $\F$ 是 $X$ 上的模层，\ednote{询问 A. G.：这里的模是否始终指凝聚或拟凝聚模层。}
以 $\G_\xi$ 表示它在 $Y_\xi$ 上的逆像；原英译稿接着写，“以 $\G$ 表示它在 $H$ 上的逆像”，从而 $\G_\xi=\G\otimes_{\OO_{P^\vee}}\OO_{S'}$。\cmt{T04}{这里 $\G$ 应当定义在 $Y=X\times_PH$ 上，而非一般没有典范态射到 $X$ 的 $H$ 上。下面的图也把 $\G$ 标在 $Y$ 上，提供了同一底本内部的更正依据。}

把主要构造与记号归纳如下（$\F,\G,\G_\eta$ 分别位于 $X,Y,Y_\eta$ 上）：
\[
\begin{tikzcd}[column sep=1.6em,row sep=2.1em]
X \arrow[d,"f"'] & X\times_SP^\vee \arrow[l] \arrow[d] & Y \arrow[l,hook'] \arrow[d] & Y_\eta \arrow[l] \arrow[d]\\
P \arrow[d] & P\times_SP^\vee \arrow[l] \arrow[d] & H \arrow[l,hook'] \arrow[dl] & H_\eta \arrow[l] \arrow[d]\\
S & P^\vee \arrow[l] & & S' \arrow[ll,"\xi"]
\end{tikzcd}
\]
图中的方形与菱形都为笛卡尔图。下一节将系统研究如下情况：$S'=\Spec K$，其中 $K$ 为域，它在 $P^\vee$ 中的像是相应纤维 $P_s^\vee$ 的泛点。作基变换 $\Spec k(s)\to S$ 后，归结到 $S=\Spec k$、$k$ 为域；下一节始终作此假设。

所研究的 $X,Y_\xi$ 的多数性质具有“几何性质”，因此在基变换下不变。这使我们可不失一般性地限制到 $K$ 代数闭，或者 $K=k(\eta)$（$\eta$ 为 $P^\vee$ 的泛点、$\xi:\Spec K\to P^\vee$ 为典范态射）的情形。对于涉及 $X,Y_\xi$ 的几何问题，亦可在基变换后设 $k$ 代数闭。

术语提醒：若 $f$ 为浸入，通常称 $Y_\xi$ 为 $X$ 的\textbf{超平面截面}（相对于射影浸入 $f$ 与超平面 $H_\xi$）。没有理由不把这个术语推广到任意 $f$ 的情形。

\egasection{§2\quad 泛超平面截面的研究：局部性质}
现在 $S=\Spec k$，$k$ 为域。若 $\eta$ 是 $P^\vee$ 的点，$\xi:\Spec k(\eta)\to P^\vee$ 为典范态射，则也写 $H_\eta,Y_\eta,\G_\eta$，代替 $H_\xi,Y_\xi,\G_\xi$。在本节中，$\eta$ 始终表示 $P^\vee$ 的泛点。

\egas{2.1}{命题}
设 $X$ 不可约。则 $Y_\eta$ 不可约或为空；在非空情形下，它占优于 $X$。〔原英译接有“字迹不清，向 AG 核实”以及“$Y$ 不可约”的编者文字。〕\ednote{向 Grothendieck 询问：例如第一节中带下划线的大写字母的含义和作用是什么？}
\begin{proof}
由于 $H\to P$ 是射影丛，$Y\to X$ 也是，所以 $X$ 不可约时 $Y$ 不可约。因此，$Y$ 在 $P^\vee$ 上的泛纤维 $Y_\eta$ 不可约或为空；非空时，它的泛点就是 $Y$ 的泛点，因而位于 $X$ 的泛点之上。
\end{proof}

\egas{2.2}{命题}
设 $Z$ 为 $P$ 的一个子集。则它在 $H_\eta$ 中的逆像 $Z_\eta$ 为空，当且仅当 $Z$ 的每个点都为闭点。特别地，若 $Z$ 可构造，则 $Z_\eta=\varnothing$ 当且仅当 $Z$ 有限。
\begin{proof}
可归结到 $Z$ 只有一个点 $z$；只须证明 $H_\eta$ 在 $P$ 中的像恰好由非闭点构成。令 $X$ 为 $z$ 的闭包，使用（2.1），只须证明 $Z_\eta=\varnothing$ 当且仅当闭子概形 $X\subset P$ 有限。

把 $X$ 换成 $X_{k(\eta)}\hookrightarrow P_{k(\eta)}$ 后，“必要”方向来自下面的事实——这里应给出参考，而且值得作为引理重述：若 $Y$ 是 $X$ 的某个超平面截面，且该截面为空，则 $X$ 有限；事实上，$X\subset P\setminus H$ 同时仿射和射影……〔底本此处在句中又写 $Y_\eta=\varnothing$，记号未统一。〕

“充分”方向是直接的。例如，注意 $Y$ 是 $X$ 上相对维数为 $n-1$ 的射影丛（$n$ 为 $P,P^\vee$ 在 $S$ 上的相对维数）。当 $X$ 在 $k$ 上有限时，$Y$ 的绝对维数为 $n-1$，小于 $\dim P^\vee=n$，所以 $Y\to P^\vee$ 不占优，泛纤维 $Y_\eta$ 为空。
\end{proof}

\egas{2.3}{推论}
设 $f:X\to P$ 为有限型态射，$Z$ 为 $X$ 的可构造子集。它在 $Y_\eta$ 中的逆像为空，当且仅当 $f(Z)$ 有限。特别地，$Y_\eta$ 为空，当且仅当 $f(X)$ 有限。

\egas{2.4}{推论〔第二断言有误〕}
设 $Z,Z'$ 为 $X$ 的两个闭子集，且 $Z$ 不可约；记它们在 $Y_\eta$ 中的逆像为 $Z_\eta,Z'_\eta$。则 $Z_\eta\subset Z'_\eta$ 当且仅当 $f(Z)$ 有限或 $Z\subset Z'$。

原文还声称：$Z_\eta=Z'_\eta$ 当且仅当 $f(Z),f(Z')$ 都有限，或者 $Z=Z'$。\cmt{E03}{后一断言按现有假设不成立。可以给 $Z$ 添加一个与它不相交的闭点而不改变泛截面。详见E03；若再假设 $Z'$ 也不可约，才能用前一断言的对称应用得到所写等号判据。}
\begin{proof}[原文证明]
这是（2.3）的直接结果。因为 $f(Z\setminus(Z\cap Z'))$ 只有在 $Z\subset Z'$ 或 $f(Z)$ 有限时才可能有限：若 $Z\not\subset Z'$，则 $Z\setminus(Z\cap Z')$ 在 $Z$ 中稠密，故其像在 $f(Z)$ 中稠密；若前者有限，它因可构造而闭，因此后者也有限。
\end{proof}

\egas{2.5}{推论}
对 $X$ 的每个满足 $\dim\overline{f(X_i)}>0$ 的不可约分支 $X_i$，取其在 $Y_\eta$ 中的逆像 $Y_{i\eta}$。则 $Y_{i\eta}$ 是 $Y_\eta$ 的一个不可约分支，并由此得到上述 $X_i$ 与 $Y_\eta$ 的不可约分支之间的一一对应。
\begin{proof}
由（2.3），$Y_\eta$ 是这些 $Y_{i\eta}$ 的并，且后者均闭且非空。由（2.1），它们不可约；由（2.4）的包含判据，它们之间互不包含，故结论成立。
\end{proof}
注意，若 $\dim X_i=d_i$，则 $\dim Y_i=d_i-1$。\cmt{记号}{此句底本省略了 $\eta$。依上下文应指 $Y_{i\eta}$，不能误读为普遍关联空间中位于 $X_i$ 上的整族。}
更一般地：

\egas{2.6}{命题}
假设对每个不可约分支 $X_i$ 都有 $\dim\overline{f(X_i)}>0$，亦即 $Y_{i\eta}\ne\varnothing$；或者，$f$ 为浸入且 $\dim\overline{f(X)}>0$〔此处原英译注明字迹稍不清，待向 AG 核实〕。则
\[
 \dim Y_\eta=\dim X-1.
\]
\ednote{法文为“ou que”；我认为应当译为“where”。}
\begin{proof}
由（2.5）归结到 $X$ 不可约。按构造，$Y_\eta$ 在 $X_{k(\eta)}$ 中是某个可逆模截面的除子，这个可逆模是 $\OO_P(1)$ 的逆像。$X_{k(\eta)}$ 不可约，因为 $X$ 不可约且 $k(\eta)/k$ 是纯超越扩张；这一点本应在本节开头指出。又有 $Y_\eta\ne X_{k(\eta)}$，因为 $Y_\eta\to X$ 的像并不是整个 $X$，与忠实平坦的 $X_{k(\eta)}\to X$ 不同：由（2.3），前者的像不含 $X$ 的闭点。所以
\[
 \dim Y_\eta=\dim X_{k(\eta)}-1=\dim X-1
\]
（最后一个等式需给参考）。
\end{proof}

\egas{2.7}{命题}
设 $\F$ 为 $X$ 上拟凝聚模，因而 $\G_\eta$ 在 $Y_\eta$ 上拟凝聚。取 $\F$ 的所有满足 $\dim\overline{f(Z_i)}>0$ 的关联素循环 $Z_i$。记 $Z_{i\eta}$ 为其在 $Y_\eta$ 中的逆像。则这些 $Z_{i\eta}$ 正是 $\G_\eta$ 的全部关联素循环；它们之间的包含关系与相应 $Z_i$ 之间相同。
\begin{proof}
最后一断言包含在（2.4）的包含判据中。由于 $Y\to X$ 是射影丛，所以它平坦且纤维不可约并满足 $(S_1)$。由 EGA IV 第3段，$\F$ 在 $Y$ 上的逆像 $\G$ 的关联素循环，就是 $\F$ 的关联素循环的逆像。再限制到 $Y\to P^\vee$ 的泛纤维，即知 $\G_\eta$ 的关联素循环为上述 $Z_i$ 的非空逆像；由（2.3）得到结论。

严格说，经由 $Y$ 的这一步并非必要。可直接用 $Y_\eta\to X$ 平坦且纤维满足 $(S_1)$、并不可约的事实；其纤维甚至几何正则（即该态射正则）且几何不可约，因为它们是射影概形的局部化。这一点也可用于（2.1）的证明。
\end{proof}

\egas{2.8}{命题}
设 $\F$ 在 $X$ 上凝聚，$y\in Y_\eta$，$x$ 为其在 $X$ 中的像。设 $P(M)$ 是诺特局部环 $A$ 上有限型模 $M$ 的下列性质之一：
\begin{enumerate}[label=(\roman*)]
\item $\coprof M\leq n$〔参考待补〕；
\item $M$ 满足 $(S_k)$〔参考待补〕；
\item $M$ 是 Cohen--Macaulay 模；
\item $M$ 约化〔参考待补〕；
\item $M$ 整〔参考待补〕。
\end{enumerate}
则 $\G_{\eta,y}$ 满足 $P$，当且仅当 $\F_x$ 满足 $P$。
\begin{proof}
这直接来自第6段的结果，\ednote{澄清这一参考：是 EGA IV 吗？}
并使用 $Y_\eta\to Y$ 是正则态射，从而 $\OO_{X,x}\to\OO_{Y_\eta,y}$ 正则这一事实。\cmt{T04}{要推出所写局部环同态正则，此处应援引刚在（2.7）讨论的 $Y_\eta\to X$；底本把靶写为 $Y$，需核对并区分两条态射。}
结合（2.3），得到下列结论。
\end{proof}

\egas{2.9}{推论}
沿用（2.8）的记号，设 $Z$ 为 $X$ 上不满足 $P(\F_x)$ 的点的集合。则 $\G_\eta$ 处处满足 $P$，当且仅当 $f(Z)$ 是 $P$ 的有限子集，或者说 $\dim\overline{f(Z)}=0$。\cmt{T05}{“有限”包括空集，而按第0章约定 $\dim\varnothing=-\infty$；若要包含无坏点的情形，应写维数 $\leq0$，或单独说明 $Z$ 为空的情况。}
\begin{proof}
（2.8）说明 $h^{-1}(Z)$ 是 $\G_\eta$ 的 $P$-奇异点集；由（2.3），其为空当且仅当 $f(Z)$ 有限。（这里 $h$ 表示 $Y_\eta\to X$；我刚刚发现，在（2.8）中字母 $P$ 被使用了两次。）
\end{proof}

\egas{2.10}{推论〔原文仅有提纲〕}
“$Y_\eta$ 约化、分别地局部整的条件。”\cmt{P01}{英文底本在此只有这一句，没有完整陈述或证明。本译稿不补造原文。}

\egas{2.11}{推论}
设 $y\in Y_\eta$，$x$ 为其像。$Y_\eta$ 在 $y$ 处正则，分别地满足 $(R_k)$〔参考待补〕，分别地正规，当且仅当 $X$ 在 $x$ 处满足同一性质。设 $Z$ 为 $X$ 的非正则点集，分别地不满足 $(R_k)$ 的点集，分别地非正规点集；则 $Y_\eta$ 正则、分别地满足 $(R_k)$、分别地正规，当且仅当 $f(Z)$ 有限，即 $\dim\overline{f(Z)}=0$。\cmt{T05}{最后的等号仍需像（2.9）一样处理空集。底本中还将一次 $R_k$ 排成了 $E_k$。}
\begin{proof}
证明与（2.8）、（2.9）相同。必须给出保证 $Z$ 闭的各个参考，因为为应用（2.3），必须知道它可构造。
\end{proof}
注意，（2.10）、（2.11）完全没有讨论相应的几何性质；这里的结果是“绝对”的。现在考察几何性质。（也可以在这里另起一节。）

\subsec{几何性质〔英译编者所加标题〕}
\egas{2.12}{定理〔原证明含错误步骤〕}
设 $f:X\to P$ 非分歧。设 $y\in Y_\eta$，$x$ 为其在 $X$ 中的像。则 $X$ 在 $x$ 处对 $k$ 光滑，当且仅当 $Y_\eta$ 在 $y$ 处对 $k(\eta)$ 光滑。
\begin{proof}[原文证明]
可设 $k$ 代数闭。若 $Y_\eta$ 在 $y$ 处对 $k(\eta)$ 光滑，则它在该点正则；$Y_\eta\to X$ 平坦，所以 $X$ 在 $x$ 处正则〔参考待补〕。由于 $k$ 代数闭，因而完美，$X$ 在 $x$ 处对 $k$ 光滑〔参考待补〕。

反过来，可将 $X$ 换成 $x$ 的开邻域而设 $X$ 光滑；原文进而说，由光滑性的 Jacobian 判据，可以在 $P$ 的某个开子集 $U$ 中用 $p$ 个方程表示为
\[
 X=V(f_1,\ldots,f_p),
\]
其中 $df_1,\ldots,df_p$ 处处线性无关。\cmt{G01}{非分歧态射不自动在 Zariski 局部成为浸入；此处还需适当的 étale 局部化或另作论证。详见G01。}
在 $P^\vee$ 中引入仿射坐标 $S_1,\ldots,S_n$，在 $x$ 的邻域中引入 $P$ 的仿射坐标 $T_1,\ldots,T_n$（选取不包含 $x$ 的无穷远超平面 $H_\infty$），则 $Y_\eta\hookrightarrow U_{k(\eta)}$ 为
\[
 V\left(f_1,\ldots,f_p,\sum_iS_iT_i-1\right).
\]
只须验证 $f_1,\ldots,f_p,\sum_iS_iT_i-1$ 相对于 $k(\eta)$ 的微分线性无关。这些微分是 $\Omega^1_{U/k}\otimes_kk(\eta)$ 的截面
\[
 df_1,\ldots,df_p,\sum_iS_i\,dT_i.
\]
由于 $df_i$ 在 $U$ 的每一点线性无关，而且 $dT_i$ 处处构成 $\Omega^1_{U/S}$ 的基，因而构成生成系，原文断言这些微分在 $U_{k(\eta)}$ 的每一点都线性无关，至少在 $p\leq n$、即
\[
 E=\Omega^1_{U/k}\Big/\sum_{1\leq i\leq n}\OO_Udf_i\ne0
\]
时如此。\cmt{T06}{底本在商的求和范围写到 $n$，而已列出的方程只有 $f_1,\ldots,f_p$；此外，商的秩为 $n-p$，非零对应 $p<n$，不是 $p\leq n$。}
这是关于非零局部自由模 $E$ 的一族生成元 $a_i$（$1\leq i\leq n$）的“小引理”：截面 $\sum_iS_ia_i$ 在 $E\otimes_kk(\eta)$ 中处处不消失。\cmt{E04}{这条辅助断言为假：在 $\mathbf A^1_k$ 的平凡线丛上取生成元 $1,t$，$S_1+S_2t$ 在 $t=-S_1/S_2$ 处为零。此反例否定这一步，不否定定理本身。详见E04及G01。}
另一方面，$p=n$ 时是平凡情形，因为 $Y_\eta=\varnothing$。
\end{proof}

\egas{2.13}{推论}
设 $Z$ 为 $X$ 不在 $k$ 上光滑的点集。则 $Y_\eta$ 对 $k(\eta)$ 光滑，当且仅当 $Z$ 有限。特别地，若 $X$ 光滑，则 $Y_\eta$ 也光滑。
\begin{proof}
由（2.12）、（2.3）得出。更一般地，有以下结果。\cmt{假设继承}{这里继承（2.12）的非分歧假设；非分歧有限型态射拟有限，因此有限的 $f(Z)$ 与有限的 $Z$ 在此可联系起来。不能把这一推论脱离其前提独立引用。}
\end{proof}

\egas{2.14}{定理}
设 $y\in Y_\eta$，$x$ 为其在 $X$ 中的像。设 $P(A,K)$ 为域 $K$ 上诺特局部代数 $A$ 的下列性质之一：
\begin{enumerate}[label=(\roman*)]
\item $A$ 在 $K$ 上几何正则；
\item $A$ 在 $K$ 上几何满足 $(R_k)$；
\item $A$ 在 $K$ 上可分；
\item $A$ 在 $K$ 上几何正规。
\end{enumerate}
原英译稿写出的结论是\cmt{T07}{右端底域应核为 $k(\eta)$，不是 $k(y)$。局部环有到剩余域 $k(y)$ 的商映射，却一般没有相反方向的典范系数域嵌入；因而底本右端甚至没有给出所需的代数结构。}
\[
 P(\OO_{X,x},k)\Longleftrightarrow P(\OO_{Y_\eta,y},k(y)).
\]

\ednote{旁注：“$X$ 非分歧，或 $k$ 的特征为0。”〔这里按与上下文相符的意义，应理解为态射 $f$ 非分歧。〕}
\begin{proof}[原文证明]
考虑第6段的结果后，（iii）、（iv）由（ii）及（2.8，（ii））得出。另一方面，（i）已在（2.12）中证明；只须由（i）推出（ii）。令 $Z$ 为 $X$ 不光滑于 $k$ 的点集，则由（2.12），其逆像 $Z_\eta$ 正是 $Y_\eta$ 不光滑于 $k(\eta)$ 的点集。由平坦性，$Z$ 在 $X$ 中于 $x$ 处的余维数等于 $Z_\eta$ 在 $Y_\eta$ 中于 $y$ 处的余维数〔第6段的参考待补〕。因而其中一个大于 $k$ 当且仅当另一个也大于 $k$，证明完成。
\end{proof}

\egas{2.15}{推论}
设 $Z$ 为 $X$ 不在 $k$ 上光滑的点集，分别地不在 $k$ 上几何满足 $(R_k)$、不可分、非正规的点集。则 $Y_\eta$ 对 $k(\eta)$ 光滑，分别地几何满足 $(R_k)$、可分、几何正规，当且仅当 $Z$ 有限。\cmt{P04}{本句的最后一项左侧只写“非正规”，右侧却写“几何正规”，应核对是否漏掉“几何”；又若采用脚注允许的“特征0而 $f$ 任意”的版本，不能把“$f(Z)$ 有限”不加条件地换成“$Z$ 有限”。详见E05及P04。}

从编写角度看，（2.14）、（2.15）包含（2.12）、（2.13），后两者可以不再独立陈述；另一方面，推论在实际使用中比定理更重要，可以把定理放在一个预备命题或引理里，使推论更为突出。对于性质（iii），还可给出一个变体：

\egas{2.16}{推论}
（仍设 $f$ 为浸入，且 $\F$ 凝聚。）在（2.7）的条件下，$Z_i$ 不是嵌入关联分支，当且仅当 $Z_{i\eta}$ 不是嵌入关联分支。\cmt{用语}{底本写“should not be immersed”；在关联素循环语境中，这里应理解为“不是嵌入分支”，不是态射的浸入性质。由于未核对法文手稿，此处明确标示这一译校理解。}
此时，$\F$ 在 $Z_i$ 处相对于 $k$ 的根重数，等于 $\G_\eta$ 在 $Z_{i\eta}$ 处相对于 $k(\eta)$ 的根重数。
\begin{proof}[原文证明]
第一断言包含在（2.7）的最后一断言中。为证第二断言，因为 $Y_\eta\to X$ 平坦，所以 $\F\mapsto\G_\eta$ 正合；限制到 $Z_i$ 泛点的邻域并使用 dévissage（逐层分解），归结为 $\F=\OO_{Z_i}$，从而可设 $Z_i=X$。接着原英译写道：“也可先假设 $X$ 在 $k$ 上分离，归结到 $k$ 代数闭。”\ednote{无法理解。}\cmt{P01}{此句在所上传英译中语法和数学衔接均未完整；这里不擅自补成一个有效的下降论证。}
于是该断言包含在（2.15，（iii））中。像通常一样，再得到下面的推论。
\end{proof}

\egas{2.17}{推论}
设 $Z$ 为 $X$ 上 $\F$ 不对 $k$ 可分的点集〔参考待补〕。则 $\G_\eta$ 对 $k(\eta)$ 可分，当且仅当 $Z$ 有限。特别地，若 $\F$ 对 $k$ 可分，则 $\G_\eta$ 对 $k(\eta)$ 可分。

\egas{2.18}{备注}
若舍弃 $f$ 是浸入的假设，关键结果（2.12）及其推论（2.13）、（2.17）便不再成立；例如可取 $X$ 为 $P$ 的一个纯不可分覆盖。但是，当 $k$ 的特征为0时，无须假设 $f$ 为浸入，（2.12）、（2.17）也成立。
\begin{proof}[原文论证]
只须对（2.12）验证；它由（2.11）以及特征0的代数预概形上“光滑等于正则”的事实得出。
\end{proof}

% 接续：第一篇3.1；范围已在卷首与进度文件标明。
